\section*{Hoofdstuk \ref{ch:g2:where-animals-live} --- Waar dieren wonen}

\begin{solution}{exo:g2:where-animals-live:1}
Het soort plek waar een \omterm{def:g1:animals-around-us:animal}{dier} van nature woont. Het vindt er zijn
voedsel, water, schuilplaats en een veilig plekje voor zijn jongen.
\end{solution}

\begin{solution}{exo:g2:where-animals-live:2}
Vijver: kikker, libel (of stekelbaars, poelslak). Bos: eekhoorn, specht
(of hert, bosmier). Wei: konijn, sprinkhaan (of vlinder, veldmuis).
\end{solution}

\begin{solution}{exo:g2:where-animals-live:3}
Stekelbaars: de vijver. Specht: het bos. Sprinkhaan: de wei. Pissebed:
onder de stenen van een oude muur (vochtige verborgen hoekjes).
\end{solution}

\begin{solution}{exo:g2:where-animals-live:4}
De kikker: zwemvliezen die het water wegduwen. De eekhoorn: grijpende
klauwen en een grote balanceerstaart.
\end{solution}

\begin{solution}{exo:g2:where-animals-live:5}
Brede graafschoppen als handen; een \omterm{def:g1:animals-around-us:covering}{vacht} die beide kanten op glad
ligt zodat de aarde hem nooit in de war maakt; \omterm{def:g1:the-five-senses:sense}{ogen} die hij in het
donker nauwelijks nodig heeft.
\end{solution}

\begin{solution}{exo:g2:where-animals-live:6}
Specht: hoog op de stammen. Eekhoorn: in de takken. Hert: tussen de
struiken. Bosmier: op de grond, in zijn opgehoopte nest.
\end{solution}

\begin{solution}{exo:g2:where-animals-live:7}
Ze verdwijnen met de vijver --- zonder het \omterm{def:g2:where-animals-live:habitat}{leefgebied} verliezen ze in
één klap voedsel, schuilplaats en voortplantingsplek. \omterm{def:g1:animals-around-us:animal}{Dieren} beschermen
begint dus bij het beschermen van de plekken waar ze wonen.
\end{solution}

\begin{solution}{exo:g2:where-animals-live:8}
Open antwoord. Zelfs onder één omgedraaide steen zitten vaak
pissebedden, mieren, een kever, soms een duizendpoot --- kleine
\omterm{def:g2:where-animals-live:habitat}{leefgebieden} zijn verrassend vol.
\end{solution}

\begin{solution}{exo:g2:where-animals-live:9}
Zwemvliezen (peddelen), \omterm{def:g1:animals-around-us:covering}{veren} die het water laten aflopen, een zeefsnavel:
alle drie wijzen naar water --- een vijver- of meerleefgebied. Het \omterm{def:g1:animals-around-us:animal}{dier}
is heel waarschijnlijk een eend.
\end{solution}

\begin{solution}{exo:g2:where-animals-live:10}
De uitrusting van de kikker --- zwemvliezen, een vochtige \omterm{def:g1:the-five-senses:sense}{huid}, een
lichaam gemaakt voor water en natte oevers --- is nutteloos in een
droge bak, en zijn \omterm{def:g2:where-animals-live:habitat}{leefgebied} gaf hem meer dan voedsel: water voor zijn
\omterm{def:g1:the-five-senses:sense}{huid} en zijn eitjes, schuilplaatsen, andere kikkers. Vriendelijk zijn
voor een wild \omterm{def:g1:animals-around-us:animal}{dier} is het in zijn \omterm{def:g2:where-animals-live:habitat}{leefgebied} laten --- of het daarheen
terugbrengen.
\end{solution}
